\documentclass[11pt]{amsart}

\usepackage[T1]{fontenc}
\usepackage{lmodern}
\usepackage{microtype}
\usepackage{amsmath,amssymb,amsthm}
\usepackage[hidelinks]{hyperref}

\hypersetup{
  pdftitle={Nine Mutually Orthogonal Cyclic 9-Cycle Systems of K19},
  pdfsubject={Mutually orthogonal cyclic cycle systems},
  pdfkeywords={orthogonal cycle systems, cyclic cycle systems, cycle decompositions}
}

\newtheorem{theorem}{Theorem}
\newtheorem{lemma}[theorem]{Lemma}

\newcommand{\Z}{\mathbb Z}

\title{Nine Mutually Orthogonal Cyclic $9$-Cycle Systems of $K_{19}$}
\date{}

\subjclass[2020]{05B30, 05C51}
\keywords{mutually orthogonal cycle systems, cyclic cycle systems,
cycle decompositions}

\begin{document}

\begin{abstract}
Burgess, Cavenagh, and Pike reported eight mutually orthogonal cyclic
$9$-cycle systems of $K_{19}$ and stated that they did not know whether a
larger family exists. We give an explicit family of nine. Consequently,
\[
9\leq \mu'(9,19)\leq 16.
\]
\end{abstract}

\maketitle

An $\ell$-cycle system of $K_n$ is a partition of $E(K_n)$ into
$\ell$-cycles. Two such systems are \emph{orthogonal} if every cycle in one
shares at most one edge with every cycle in the other. A system on vertex
set $\Z_n$ is \emph{cyclic} if it is invariant under translation. Let
$\mu'(\ell,n)$ be the maximum size of a pairwise orthogonal family of cyclic
$\ell$-cycle systems of $K_n$.

Burgess, Cavenagh, and Pike found a family of size eight for
$(\ell,n)=(9,19)$ and left open whether a larger family exists
\cite[Appendix, Table~1]{BCP}. We answer this question affirmatively.

\begin{theorem}\label{thm:main}
There exist nine mutually orthogonal cyclic $9$-cycle systems of $K_{19}$.
In particular, $\mu'(9,19)\geq 9$.
\end{theorem}

For the proof, let $C$ be a cycle in $\Z_{19}$ and, for $t\in\Z_{19}$, let
$C+t$ be obtained by adding $t$ to every vertex. Set
\[
\mathcal O(C)=\{C+t:t\in\Z_{19}\}.
\]
The \emph{difference} of an edge $\{x,y\}$ is the unique
$d\in\{1,\ldots,9\}$ such that $y-x\equiv\pm d\pmod{19}$.

\begin{lemma}\label{lem:criterion}
Let $C$ and $D$ be $9$-cycles, each containing exactly one edge of every
difference $d\in\{1,\ldots,9\}$. Orient each difference-$d$ edge in the
positive-$d$ direction, and let $a_d$ and $b_d$ be its tail in $C$ and $D$,
respectively. Then $\mathcal O(C)$ and $\mathcal O(D)$ are orthogonal if and
only if the residues
\[
a_d-b_d,\qquad 1\leq d\leq 9,
\]
are pairwise distinct modulo $19$.
\end{lemma}

\begin{proof}
For $x,y\in\Z_{19}$, the difference-$d$ edges of $C+x$ and $D+y$ coincide
if and only if
\[
a_d+x=b_d+y,
\]
or equivalently $a_d-b_d=y-x$. Thus the number of common edges of $C+x$
and $D+y$ is the multiplicity of $y-x$ among the nine residues
$a_d-b_d$. The result follows.
\end{proof}

\begin{proof}[Proof of Theorem~\ref{thm:main}]
All arithmetic is in $\Z_{19}$. Let
\[
C=(0,1,3,18,5,13,4,9,12)
\qquad\text{and}\qquad
\mathcal F_0=\mathcal O(C).
\]
Including the closing edge, the signed consecutive differences of $C$ are
\[
(1,2,-4,6,8,-9,5,3,7),
\]
so $C$ contains exactly one edge of every difference $1,\ldots,9$. For each
$d$, translating its difference-$d$ edge gives all $19$ edges of that
difference class exactly once. Since $9\cdot19=\binom{19}{2}$,
$\mathcal F_0$ is a cyclic $9$-cycle system of $K_{19}$.

The element $4\in\Z_{19}^{\times}$ has order $9$: one has
$4^9\equiv1\pmod{19}$, whereas $4\not\equiv1$ and
$4^3\equiv7\not\equiv1\pmod{19}$. For $0\leq r\leq8$, define
\[
\mathcal F_r=4^r\mathcal F_0.
\]
Multiplication by $4^r$ permutes $\Z_{19}$, and hence
\[
4^r\mathcal F_0
 =\{4^rC+4^rt:t\in\Z_{19}\}
 =\mathcal O(4^rC).
\]
Thus every $\mathcal F_r$ is a cyclic $9$-cycle system.

For $\mathcal F_0$, the tails of the positively oriented difference-$d$
edges, indexed by $d=1,\ldots,9$, are
\[
a^{(0)}=(0,1,9,18,4,18,12,5,4).
\]
Let $a^{(r)}$ be the corresponding tail vector for $4^rC$. Direct
calculation gives
\[
\begin{array}{c|rrrrrrrrr}
 r\backslash d&1&2&3&4&5&6&7&8&9\\ \hline
 1& 3& 6&16&18& 8& 4& 2&1&13\\
 2&16& 1&12&11& 5& 8& 9&6& 2\\
 3& 4&14& 3& 6& 2& 9&12&5&11\\
 4&11&12& 2&15& 6& 5&14&0& 4
\end{array}
\qquad
\bigl(a^{(0)}_d-a^{(r)}_d\bigr)_{d=1}^{9}.
\]
Each row contains nine distinct residues modulo $19$. By
Lemma~\ref{lem:criterion},
\[
\mathcal F_0\perp\mathcal F_r
\qquad (1\leq r\leq4).
\]

Multiplication by a nonzero residue preserves edge intersections. Hence any
pair $\mathcal F_i,\mathcal F_j$ is carried, after multiplication by
$4^{-i}$, to $\mathcal F_0,\mathcal F_{j-i}$, with exponents modulo $9$.
Orthogonality is symmetric, so the cases $r$ and $9-r$ are equivalent. The
four checks above therefore cover every pair, and
$\mathcal F_0,\ldots,\mathcal F_8$ are mutually orthogonal.
\end{proof}

By Theorem~\ref{thm:main} and the general bound
$\mu'(\ell,n)\leq n-3$ from \cite[Lemma~6.4]{BCP},
\[
9\leq\mu'(9,19)\leq16.
\]
The construction does not rule out a family of size ten or larger.

\begin{thebibliography}{1}

\bibitem{BCP}
A.~C. Burgess, N.~J. Cavenagh, and D.~A. Pike,
\emph{Mutually orthogonal cycle systems},
Ars Math. Contemp. \textbf{23} (2023), no.~2, Paper No.~P2.05, 20 pp.,
\href{https://doi.org/10.26493/1855-3974.2692.86d}
{doi:10.26493/1855-3974.2692.86d}.

\end{thebibliography}

\end{document}